%% ma: L11-B09-0004
%% lop: 11
%% chuong: 3
%% bai: 9
%% dang: 11.9.1
%% loai: TN4
%% muc: VD
%% dap_an: "D"
%% nguon: "chuyen-de-cac-so-dac-trung-cua-mau-so-lieu-khong-ghep-nhom-toan-10-chuong-trinh-moi.pdf (GV. Phan Nhật Linh) (P02), Bài 02 Dạng 1 Phần I Câu 33"
%% kien_thuc: "Lớp 11 Bài 9 (số trung bình của mẫu số liệu ghép nhóm)"
%% trang_thai: chinh-thuc
%% ghi_chu: "Mẫu số liệu ghép nhóm (vượt Lớp 10 Bài 13), đã chuyển sang Lớp 11 Bài 9; đề gốc hỏi 'trung bình số người' chênh lệch nhưng thực chất là tuổi trung bình, đã sửa lời dẫn. Tự duyệt theo yêu cầu giáo viên 22/10/2026."
\begin{ex}
Số người cấp cứu ở một bệnh viện trong hai ngày thứ Hai và thứ Sáu được cho trong bảng tần số ghép nhóm sau (nhóm theo độ tuổi):
\begin{center}
\begin{tabular}{|c|c|c|c|c|c|c|c|c|}\hline
Tuổi & $[4;7]$ & $[8;11]$ & $[12;15]$ & $[16;19]$ & $[20;23]$ & $[24;27]$ & $[28;31]$ & Cộng\\\hline
Thứ Hai & $1$ & $4$ & $15$ & $26$ & $16$ & $7$ & $3$ & $72$\\\hline
Thứ Sáu & $1$ & $4$ & $21$ & $22$ & $13$ & $3$ & $1$ & $65$\\\hline
\end{tabular}
\end{center}
Tuổi trung bình của những người cấp cứu trong ngày thứ Hai và ngày thứ Sáu (làm tròn từng giá trị đến hàng đơn vị) chênh lệch nhau bao nhiêu tuổi?
\choice{$4$}{$3$}{$2$}{\True $1$}
\loigiai{Giá trị đại diện các nhóm: $5{,}5$; $9{,}5$; $13{,}5$; $17{,}5$; $21{,}5$; $25{,}5$; $29{,}5$.

Thứ Hai: $\overline x_1=\dfrac{1312}{72}\approx18$. Thứ Sáu: $\overline x_2=\dfrac{1097{,}5}{65}\approx17$.

Chênh lệch $18-17=1$ (tuổi). Chọn D.}
\end{ex}
